% --------------------------------------------------------------------------------------------
% 										Capitulo 1
% --------------------------------------------------------------------------------------------

\chapter{Conclusões}

Esta dissertação apresentou o desenvolvimento de uma pipeline modular para simulação de observações radioastronômicas do tipo Single Dish, integrando a modelagem do céu, da resposta instrumental, da geração de Time Ordered Data (TOD) e da reconstrução de mapas em uma única infraestrutura computacional. Ao longo do trabalho foram reunidos os fundamentos físicos, matemáticos e computacionais necessários para descrever, de forma integrada, a relação entre o céu observado, a resposta instrumental, a geração de \textit{Time Ordered Data} (TOD) e a reconstrução de mapas.

Uma das principais contribuições deste trabalho consiste na modelagem integrada da cadeia observacional no regime \textit{Single Dish}. A estrutura proposta estabelece uma correspondência entre os processos físicos envolvidos na observação e sua representação computacional, conectando, em uma mesma modelagem, a emissão astrofísica, a resposta do instrumento, a aquisição dos dados temporais e a etapa de reconstrução de mapas. Essa abordagem fornece uma visão sistêmica do processo observacional e estabelece uma base conceitual consistente para o desenvolvimento de simulações voltadas à investigação de efeitos instrumentais e metodologias de processamento de dados.

Outra contribuição relevante é a arquitetura desenvolvida para a pipeline de simulação. A organização do software em módulos especializados, baseada em princípios de Programação Orientada a Objetos e \textit{Domain-Driven Design}, permitiu separar as diferentes responsabilidades do sistema e representar, de forma independente, os principais elementos físicos e computacionais envolvidos na simulação. Essa estrutura favorece a reutilização e manutenção do código, além da incorporação de novos modelos instrumentais ou observacionais sem alterações significativas na arquitetura existente.

Os módulos implementados estabelecem uma infraestrutura computacional capaz de suportar as diferentes etapas do fluxo de simulação, desde a geração de mapas sintéticos do céu até a produção de TOD e sua posterior reconstrução em mapas. Dessa forma, a pipeline constitui uma base para estudos relacionados à caracterização de instrumentos, avaliação de estratégias observacionais, desenvolvimento de algoritmos de processamento e geração de conjuntos de dados sintéticos.

Como continuidade desta pesquisa, a próxima etapa será dedicada à realização dos experimentos numéricos e à validação  dos modelos implementados. Nessa fase a  pipeline séra validada por meio de diferentes procedimentos, incluindo experimentos com céus sintéticos de propriedades conhecidas, comparação dos resultados com modelos teóricos, verificação da reprodução de características observacionais esperadas e comparação com ferramentas independentes. Isso permitirá avaliar sua capacidade de reproduzir características observacionais relevantes e investigar o impacto dos diferentes componentes físicos e instrumentais incorporados ao modelo.


